\documentclass[12pt, xcolor=dvipsnames,svgnames,x11names]{article}
\usepackage{xcolor}
\usepackage{xspace}
\usepackage{url}
\usepackage[colorlinks=true,bookmarks=false,linkcolor=black,urlcolor=blue,citecolor=black]{hyperref}
\usepackage{fancyhdr}
\usepackage[top=1in,bottom=1.2in,left=1in,right=1in]{geometry}
\usepackage{multicol}
\usepackage{pgfplots}
\usepackage{tikz}
\usetikzlibrary{circuits.logic.US, patterns}
\usepackage{mathpazo}
\usepackage{biolinum} 
\usepackage[all]{tcolorbox}
\usepackage{derivative}
\usepackage{../Lectures/rahul_math}
\renewcommand{\familydefault}{\sfdefault}

% --- Custom Color Palette from Image ---
\definecolor{headerRed}{HTML}{A3403D} % Dark red from image top
\definecolor{patternBg}{HTML}{F2D9D7} % Light pinkish bg from image bottom
\definecolor{binaryText}{HTML}{D1A7A4} % Muted color for binary digits

% --- Background Pattern Setup ---
\usepackage[pages=all]{background}
\backgroundsetup{
scale=1,
color=black,
opacity=1,
angle=0,
contents={%
  \begin{tikzpicture}[remember picture,overlay]
    % Top Header Bar
    \fill[headerRed] (current page.north west) rectangle ([yshift=-1.2cm]current page.north east);
    % Bottom Binary Background
    \fill[patternBg] (current page.south west) rectangle ([yshift=3.5cm]current page.south east);
    % Decorative Line (from image)
    \draw[headerRed, line width=2pt] ([xshift=1in, yshift=-2.5cm]current page.north west) -- ([xshift=8in, yshift=-2.5cm]current page.north west);
    % Binary Text Overlay (Sample)
    \node[anchor=south, binaryText, font=\ttfamily\scriptsize, inner sep=0pt, yshift=1.5cm] at (current page.south) {
        \begin{tabular}{c@{\hspace{0.5em}}c@{\hspace{0.5em}}c@{\hspace{0.5em}}c@{\hspace{0.5em}}c@{\hspace{0.5em}}c@{\hspace{0.5em}}c}
        0 0 0 0 1 1 1 1 0 0 0 0 1 1 0 1 1 1 0 0 0 1 \\
        1 0 0 1 0 1 0 0 0 1 0 0 1 1 0 0 1 0 1 1 \\
        0 0 1 1 0 0 1 0 1 1 1 0 1 1 0 1 0 1 0 0 0 \\
        \end{tabular}
    };
  \end{tikzpicture}
}
}

% --- Header Styling ---
\makeatletter
\renewcommand{\maketitle}{%
    \vspace*{0.2cm}
    \begin{center}
        \begin{tcolorbox}[
            enhanced,
            colback=white,
            colframe=headerRed,
            arc=0mm,
            title={\textcolor{white}{\bfseries \@title}},
            coltitle=white,
            fonttitle=\bfseries\large,
            attach boxed title to top left={xshift=10pt, yshift=-\tcboxedtitleheight/2},
            boxed title style={sharp corners, colback=headerRed, frame hidden}
        ]
        \large
        \vspace{10pt}
        \noindent \textbf{Student Name:} \underline{\hspace{7cm}} \hfill \textbf{Points:} 30
        
        \vspace{10pt}
        \noindent \textbf{A Number:} \underline{\hspace{11cm}}
        \end{tcolorbox}
    \end{center}
    \vspace{10pt}
}
\makeatother

\newcommand{\hw}{Classwork 01// Matrix Calculus: Jacobian}
\newcommand{\duedate}{Jan 22, 2026}

\pagestyle{fancy}
\fancyhf{}
\lhead{\small CPE 487/587, Spring 2026}
\rhead{\small \textit{\hw}}
\cfoot{\thepage}
\renewcommand{\headrulewidth}{0pt}

\begin{document}

\title{\hw}
\maketitle


\begin{enumerate}
\item Derive the Jacobian of identify function $\fbf(\xbf) = \xbf $, i.e. $f_i(\xbf) = x_i$ that has $n$ functions and each function has $n$ parameters held in a single vector $\xbf$. \hfill \textbf{(10 Points)}

\item Under what conditions, are the off-diagonal elements of a Jacobian zero? \hfill \textbf{(10 Points)}

\item Compute gradient of $y = sum(\xbf)$, i.e. $\nabla y $, considering the function inside the summation as $f_i(\xbf) =  x_i$. Is it a vertical vector or a horizontal vector as per convention used in the Lecture notes? \hfill \textbf{(10 Points)}

\end{enumerate}


\section*{Solution}

\begin{enumerate}
\item 

Consider $\ybf = \fbf(\xbf)$ for shorthand notation. As the function given in the problem is an identity function, so

\begin{align*}
y_1 = f_1(\xbf) = x_1\\
y_2 = f_2(\xbf) = x_2\\
\vdots\\
\end{align*}

So we can write Jacobian as

\begin{align*}
\pdv{\ybf}{\xbf} = \begin{bmatrix}
\pdv{f_1(\xbf)}{\xbf} \\\pdv{f_2(\xbf)}{\xbf}\\ \vdots \\ \pdv{f_m(\xbf)}{\xbf}
\end{bmatrix} & = \begin{bmatrix}
\pdv{f_1(\xbf)}{x_1} & \pdv{f_1(\xbf)}{x_2} & \cdots & \pdv{f_1(\xbf)}{x_n}\\
\pdv{f_2(\xbf)}{x_1} & \pdv{f_2(\xbf)}{x_2} & \cdots & \pdv{f_2(\xbf)}{x_n}\\
& & \cdots & \\
\pdv{f_m(\xbf)}{x_1} & \pdv{f_m(\xbf)}{x_2} & \cdots & \pdv{f_m(\xbf)}{x_n}
\end{bmatrix}
\end{align*}





since $\pdv{x_i}{x_j} = 0$ for $j \neq i$ because $f_1$ only has $x_1$ term, $f_2$ only has $x_2$ term and so on. Hence, 


\begin{align*}
\pdv{\ybf}{\xbf} = \begin{bmatrix}
\pdv{x_1}{x_1} & 0  & \cdots & 0\\
 0 & \pdv{x_2}{x_2}  & \cdots & 0\\
& & \cdots & \\
 0 & 0 & \cdots & \pdv{x_n}{x_n} 
\end{bmatrix} =  \begin{bmatrix}
1 & 0  & \cdots & 0\\
 0 & 1 & \cdots & 0\\
& & \cdots & \\
 0 & 0 & \cdots & 1
\end{bmatrix} = I
\end{align*}
$I$ is the identity matrix.


\item When scalar functions are depenedent on one variable (feature) only such that $f_1$ only has $x_1$ term, $f_2$ only has $x_2$ term and so on, then the off diagonal elements of a Jacobian will be zero.

\item $ y = sum( \xbf)$, and the function inside the summation is just $f_i(\xbf) = x_i$ and then the gradient is 

\begin{align*}
\nabla y & = \begin{bmatrix}
\sum_i \pdv{f_i(\xbf)}{x_1} & \sum_i \pdv{f_i(\xbf)}{x_2} & \cdots & \sum_i \pdv{f_i(\xbf)}{x_n}
\end{bmatrix} =  \begin{bmatrix}
\sum_i \pdv{x_i}{x_1} & \sum_i \pdv{x_i}{x_2} & \cdots & \sum_i \pdv{x_i}{x_n}
\end{bmatrix}\\
& = \begin{bmatrix}
\pdv{x_1}{x_1} & \pdv{x_2}{x_2} & \cdots & \pdv{x_n}{x_n}
\end{bmatrix} = \begin{bmatrix}
1 & 1 & \cdots & 1] 
\end{bmatrix} = \vec{1}^\top
\end{align*}

which is a horizontal vector.

\end{enumerate}




\end{document}